\documentclass[conference]{IEEEtran}
\IEEEoverridecommandlockouts
% The preceding line is only needed to identify funding in the first footnote. If that is unneeded, please comment it out.
\usepackage{cite}
\usepackage{amsmath,amssymb,amsfonts}
\usepackage[hidelinks]{hyperref}
\usepackage{algorithmic}
\usepackage{graphicx}
\usepackage{textcomp}
\usepackage{xcolor}

\usepackage{booktabs}
\usepackage{multirow}

\def\BibTeX{{\rm B\kern-.05em{\sc i\kern-.025em b}\kern-.08em
    T\kern-.1667em\lower.7ex\hbox{E}\kern-.125emX}}
\begin{document}

\title{Joint Knowledge Distillation and Tensor-Based Compression of Deep Neural Networks}

\author{\IEEEauthorblockN{1\textsuperscript{st} Enis Furkan Kırmızı}
\IEEEauthorblockA{\textit{Artificial Intelligence \& Data Engineering} \\
\textit{Istanbul Technical University}\\
Istanbul, Türkiye \\
kirmizie23@itu.edu.tr}
}

\maketitle




%%%%% ABSTRACT %%%%%




\begin{abstract}
Deploying deep neural networks on resource-constrained devices necessitates aggressive model compression; however, traditional tensor decomposition often induces severe accuracy degradation. While Knowledge Distillation (KD) mitigates this, co-distilling architecturally heterogeneous compressed models remains fundamentally challenging due to dimensional mismatches. We introduce a novel Heterogeneous Structural Co-Distillation framework using Batch-Subspace Alignment with Attention Transfer (BSAT). Using PURSUhInT for non-redundant hint-point selection, our method simultaneously trains CP-decomposed and Tucker-decomposed student networks under a shared teacher. To enforce a shared semantic latent space, we propose a dimension-agnostic BSAT loss that aligns batch-subspaces using sign-invariant soft spectral embeddings derived from Gram-matrix eigendecomposition, eliminating the need for trainable regressors. Extensive experiments demonstrate strong efficiency-accuracy trade-offs; notably, our framework compresses a ResNet-18 by 16× in parameters and 13.2× in FLOPs while maintaining 80.10\% Top-1 accuracy on ImageNet-100.
\end{abstract}



%%%%% KEYWORDS %%%%%




\begin{IEEEkeywords}
Knowledge Distillation, Model Compression, Tensor Decomposition, Heterogeneous Co-Distillation, Batch-Subspace Alignment.
\end{IEEEkeywords}




%%%%% INTRODUCTION %%%%%




\section{Introduction}
\label{sec:intro}

\par
Modern deep neural networks (DNNs) achieve state-of-the-art perceptual
accuracy through aggressive architectural scaling, yet this scaling
precipitates a critical hardware bottleneck: the memory wall, where the cost
of fetching model parameters from DRAM eclipses available compute bandwidth
\cite{gholami2024ai}. The resulting latency and energy overhead renders
overparameterized models fundamentally unsuitable for resource-constrained
deployments such as autonomous drone navigation and mobile healthcare.
Aggressive model compression is therefore an engineering imperative for
practical edge AI.

\par
Low-rank tensor factorization—specifically CP and Tucker decomposition—offers a principled route to structural compression, but the standard sequential pipeline of decomposing pre-trained weights
followed by fine-tuning frequently leads to catastrophic accuracy degradation
\cite{gusak2019automated}. Knowledge Distillation (KD)
\cite{hinton2015distilling} addresses this by transferring a teacher's learned
representations to a compact student, yet logit-level distillation alone fails
under extreme compression: the student lacks sufficient capacity to
approximate the teacher's output distribution, a phenomenon known as the
capacity gap \cite{mirzadeh2020improved}. Hint distillation mitigates this by
aligning intermediate feature maps \cite{romero2015fitnetshintsdeepnets}, but
existing methods require trainable $1\!\times\!1$ convolutional regressors to
reconcile channel mismatches and depend on rigid spatial heuristics that
transfer redundant or noisy knowledge. The PURSUhInT framework
\cite{keser2023pursuhint} improves hint selection via $R^2_\text{CCA}$
representational clustering, yet it does not address compression itself;
moreover, standard feature-alignment losses break down under the dimensional
mismatches inherent in structurally heterogeneous, tensor-decomposed
architectures.

\par
We propose a Heterogeneous Structural Co-Distillation framework that jointly
trains a CP-decomposed student $\mathcal{S}_\text{CP}$ and a Tucker-decomposed
student $\mathcal{S}_\text{Tuck}$ under a shared teacher. Rather than
factorizing pre-trained weights, both students are constructed
\emph{from scratch} before training begins, with per-layer ranks derived once
from the teacher's weight spectrum via Empirical VBMF (EVBMF)
\cite{nakajima2013global}—transferring the teacher's observed per-layer
capacity utilization directly into each student's structural blueprint. The
core distillation loss, Batch-Subspace Alignment with Attention Transfer
(BSAT), sidesteps channel-dimension mismatches entirely by aligning
sign-invariant soft spectral embeddings extracted from the $B\!\times\!B$
Gram matrix of each feature map's batch unfolding, requiring no trainable
regressor and no shared channel count. A bidirectional coupling term
additionally constrains both students to organize their batch representations
into a shared semantic subspace, using stop-gradient projectors from the peer
student as regression targets. Hint points are provided by PURSUhInT
\cite{keser2023pursuhint}, which selects the $K$ most structurally distinct
teacher layers via k-means clustering on pairwise $R^2_\text{CCA}$ distances.
On ImageNet-100, this framework compresses ResNet-18 by $16\times$ in
parameters and $13.2\times$ in FLOPs while achieving 80.10\% Top-1 accuracy,
within 1.2\% of the uncompressed baseline.

\par
The core contributions of this work are as follows:
\begin{itemize}
    \item \textbf{The BSAT Distillation Objective:} A two-term distillation
    loss combining spatial attention matching with sign-invariant soft spectral
    batch-subspace alignment via energy-adaptive Gram-matrix eigendecomposition,
    eliminating the need for trainable regressors and accommodating arbitrary
    channel and spatial mismatches.

    \item \textbf{Heterogeneous Dual-Student Co-Distillation:} A training
    framework in which structurally disparate CP-decomposed and
    Tucker-decomposed students are trained simultaneously, regularized by a
    bidirectional batch-subspace coupling term that enforces a shared semantic
    latent space without requiring compatible layer dimensions.

    \item \textbf{Teacher-Guided EVBMF Rank Selection:} Per-layer intrinsic
    ranks estimated from the teacher's fully-trained weight spectrum and
    transferred to the student's decomposed architecture via positional
    interpolation, replacing uninformative global ratios.

    \item \textbf{Non-Redundant Hint Selection via PURSUhInT:} 
    We adopt the PURSUhInT mechanism to select $K$ structurally distinct
    teacher layers via $R^2_\text{CCA}$ representational clustering, and
    demonstrate for the first time its compatibility with simultaneously
    supervising two architecturally heterogeneous compressed students.
\end{itemize}

\par
The remainder of this paper is organized as follows. Section~\ref{sec:related}
reviews related work on model compression, knowledge distillation, and
subspace alignment. Section~\ref{sec:methodology} details the proposed framework,
the BSAT loss, and the dual-student architecture. Section~\ref{sec:experiments}
presents experimental results and ablation studies on CIFAR-100 and
ImageNet-100. Section~\ref{sec:conclusion} concludes with a summary of
findings and directions for future work.




%%%%% RELATED WORK %%%%%




\section{Related Work}
\label{sec:related}


\subsection{Knowledge Distillation}
\par
Knowledge distillation (KD) \cite{hinton2015distilling} transfers the
functional behavior of an overparameterized teacher to a compact student,
offering a training-time alternative to pruning \cite{han2015deep} and
quantization \cite{jacob2018quantization}.


\subsubsection{The Evolution of Logit-Based Distillation}
\par
Classical KD uses a temperature scalar ($T>1$) on the softmax function to expose the teacher's "dark knowledge"—the secondary probability distributions between classes. Recent literature has expanded this output-driven alignment: Spherical KD (SKD) \cite{guo2021reducingteacherstudentgapspherical} normalizes logits to preserve angular geometry, Decoupled KD (DKD) \cite{zhao2022decoupled} separates target and non-target objectives to preserve critical correlations, and Weighted Soft Labels (WSL) \cite{zhou2021rethinking} dynamically calibrate loss based on sample difficulty.


\subsubsection{Structural Breakdowns: The Capacity Gap}
\par
Despite these advancements, logit distillation fails under extreme compression ratios. When a heavily factorized student attempts to distill a deep teacher, it
frequently suffers from the ``curse of the capacity gap''
\cite{mirzadeh2020improved}: the student's reduced parameterization is
insufficient to reproduce the teacher's output distribution through logit
matching alone. Logit distillation alone is therefore insufficient for extremely low-rank tensor-decomposed students, which require intermediate feature-level guidance to bridge the capacity gap.


\subsection{Feature and Attention-Based Distillation}
\par
To bypass the capacity gap, hint distillation forces the student to replicate the teacher's intermediate feature maps, converting an intractable global mapping problem into localized regression tasks.


\subsubsection{Foundational Feature Alignment Mechanisms}
\par
Methods like FitNets \cite{romero2015fitnetshintsdeepnets} train the student to regress specific middle-layer outputs before global distillation. Attention Transfer (AT) \cite{zagoruyko2016paying} shifts this alignment to spatial attention maps, teaching the student where to allocate focus rather than strictly what features to extract. Variational Information Distillation (VID) \cite{ahn2019variational} and Attention-based Feature Distillation (AFD) \cite{ji2021show} further integrate information-theoretic optimization and dynamic search mechanisms to link intermediate hints.


\subsubsection{Vulnerabilities: Trainable Regressors and Spatial Heuristics}
\par
These techniques exhibit severe architectural vulnerabilities. First, mismatched channel dimensions require the insertion of trainable $1\times1$ convolutional regressors or projection matrices. These inflate memory overhead and introduce optimization interference. Second, they depend on rigid spatial heuristics (e.g., matching the final layers of spatial blocks), assuming a fixed correspondence between teacher and student feature hierarchies that does not hold under heterogeneous compression.


\subsubsection{Systematic Layer Clustering (PURSUhInT)}
\par
To replace flawed heuristics, PURSUhInT \cite{keser2023pursuhint} quantifies semantic similarity between all teacher layers using Centered Kernel Alignment (CKA) \cite{kornblith2019similarity} and the mean squared Canonical Correlation Analysis coefficient ($R_{CCA}^2$). By applying k-means clustering to the $L$ layer representations using $R_{CCA}^2$ as the distance function, it groups redundant processing states and selects the mathematical centroid of each cluster as the optimal, non-redundant hint point. While PURSUhInT solves hint \emph{selection}, the broader class of feature alignment methods it informs remains constrained by trainable regressors and Euclidean feature losses that break down under the dimensional mismatches inherent in heterogeneous tensor-decomposed architectures.


\subsection{Tensor Decomposition for Model Compression}
\par
Low-rank tensor factorization provides a mathematically rigorous mechanism to structurally minimize the parameter space prior to inference, offering a deterministic alternative to heuristic pruning.


\subsubsection{Mechanics of CP and Tucker Decompositions}
\par
Canonical Polyadic (CP) decomposition isolates a weight tensor into a minimal sum of rank-one tensors, replacing dense convolutions with efficient spatial and channel-wise separable operations \cite{lebedev2014speeding}. However, this rigidly restricts cross-channel interactions. Tucker decomposition, conversely, preserves cross-channel correlations by utilizing a dense, low-dimensional core tensor multiplied by orthogonal factor matrices, providing independent control over input and output subspaces \cite{kim2015compression}.


\subsubsection{Structural Fragility and Uniform Factorization}
\par
Mathematical decomposition inherently disrupts pre-trained weight distributions. When fine-tuned, these networks frequently suffer from vanishing gradients and catastrophic accuracy degradation. Existing methodologies typically deploy a single factorization technique (either CP or Tucker) uniformly across the network, attempting to stabilize it through multi-stage iterative fine-tuning. Deploying a single decomposition uniformly also ignores the complementary strengths of CP and Tucker; and without distillation, multi-stage fine-tuning remains fragile under extreme compression ratios \cite{gusak2019automated}.


\subsection{Rank Estimation and Adaptive Compression}
\par
The efficacy of tensor factorization relies on determining the optimal tensor rank, which dictates the strict trade-off between compression and representational capacity.


\subsubsection{Variational Bayesian Matrix Factorization (VBMF)}
\par
Determining the optimal rank of a tensor is NP-hard in general \cite{hillar2013most}. The literature utilizes analytical techniques like Empirical Variational Bayesian Matrix Factorization (EVBMF) \cite{nakajima2013global} to dynamically denoise matrices under low-rank assumptions. EVBMF estimates the intrinsic rank of a weight matrix by applying Bayesian shrinkage to its singular values. This was successfully operationalized in pipelines like MUSCO \cite{gusak2019automated}, which estimate static ranks prior to factorization.


\subsubsection{Limitations of Static Rank Estimation}
\par
Applying EVBMF to the student's own randomly-initialized weights, or relying on uniform global rank ratios, fails to leverage the teacher's richer, fully-trained weight spectrum. The resulting ranks reflect neither the teacher's actual per-layer capacity utilization nor the heterogeneous information density across layers.


\subsection{Geometry-Aware and Subspace Alignment Methods}
\par
Because heavily factorized students cannot achieve the numerical parity required by Euclidean feature matching, geometry-aware distillation transfers the underlying mathematical structure and principal directions of variance.


\subsubsection{Relational Distillation and Spectral Filtering}
\par
Methods like Relational KD (RKD) \cite{park2019relational} and Similarity-Preserving KD (SPKD) \cite{tung2019similarity} align inter-sample angles and covariance matrices. Contrastive Representation Distillation (CRD) \cite{tian2019contrastive} demonstrates that aligning the full representation manifold—rather than individual feature vectors—preserves intrinsic task geometry more faithfully than raw feature matching. Spectral approaches such as SSKD \cite{xu2020knowledge} further show that isolating the dominant variance directions via low-rank approximation improves representation alignment under capacity constraints.


\subsubsection{The Eigenvector Sign Ambiguity Problem}
\par
Eigendecomposition ($G = U \Lambda U^T$) introduces severe numerical instability: eigenvector sign ambiguity. Both $U$ and $-U$ represent valid mathematical decompositions of the exact same subspace. Minimizing the direct distance between teacher and student eigenvectors results in arbitrary sign flipping, generating massive gradient spikes during backpropagation. Direct eigenvector alignment is therefore numerically unstable, requiring a sign-invariant representation of the learned subspace to enable clean gradient flow during distillation.


\subsection{Heterogeneous Architecture and Mutual Distillation}
\par
Collaborative, multi-agent topologies address the constraints of static hierarchies. Deep Mutual Learning (DML) \cite{zhang2018deep} demonstrated that students optimizing a peer-mimicry loss alongside supervised loss prevents overfitting and outperforms static distillation.


\subsubsection{Challenges in Cross-Architecture Distillation}
\par
Extending collaborative learning to structurally compressed, heterogeneous topologies—such as distilling from Vision Transformers to CNNs—is computationally expensive. Methods like Heterogeneous Complementary Distillation (HCD) \cite{xu2026heterogeneouscomplementarydistillation} rely heavily on complex mapping modules and sub-logit decoupled distillation to map incompatible features.


\subsubsection{Meta-Optimization and Bidirectional Coupling}
\par
To bridge low-rank compression and distillation, the Over-Parameterization Distillation Framework (OPDF) \cite{zhan2024overparameterizedstudentmodeltensor} temporarily inflates a student's parameters during training to absorb knowledge before contracting. Existing multi-student frameworks either rely on complex mapping modules for cross-architecture alignment or fail to maintain strict low-rank bounds during collaborative training.


\subsection{Summary}
\par
The reviewed literature reveals a confluence of open problems: logit distillation
collapses under extreme compression; feature-alignment methods require trainable
regressors incompatible with heterogeneous decomposed architectures; single-topology
decomposition ignores complementary factorization strengths; rank selection from
student weights misses teacher capacity signals; subspace alignment is destabilized
by eigenvector sign ambiguity; and no existing framework co-trains structurally
heterogeneous compressed students with a shared semantic constraint.
The proposed framework addresses all of these simultaneously, as detailed in Section~\ref{sec:methodology}.





%%%%% METHODS %%%%%





\section{Methodology}
\label{sec:methodology}

\subsection{Framework Overview}
\par
The proposed framework jointly compresses a pre-trained teacher network 
$\mathcal{T}$ into two structurally heterogeneous student networks—a 
CP-decomposed student $\mathcal{S}_{\mathrm{CP}}$ and a Tucker-decomposed 
student $\mathcal{S}_{\mathrm{Tuck}}$—through a unified five-stage pipeline 
(Fig.~\ref{fig:pipeline}). In the first stage, $\mathcal{T}$ is trained to 
convergence on the target dataset. In the second and third stages, the PURSUhInT 
mechanism \cite{keser2023pursuhint} extracts intermediate layer representations 
from $\mathcal{T}$, computes pairwise representational distances, and applies
k-means clustering to select $K$ non-redundant hint points
$\mathcal{H} = \{h_1, \ldots, h_K\}$—the teacher layers that carry the most
structurally distinct semantic information. In the fourth stage, both students
are constructed by decomposing the student architecture via CP and Tucker
factorization respectively, with per-layer ranks determined once before
training by running Empirical VBMF on $\mathcal{T}$'s trained weights. The
two students are then co-trained from scratch under a shared teacher,
supervised by the proposed Batch-Subspace Alignment with Attention Transfer
(BSAT) loss at each hint point and a bidirectional coupling term that
constrains both students to compress their representations into a shared
semantic latent space. In the fifth stage, both compressed models are
evaluated independently on accuracy, parameter count, FLOPs, and inference
latency. Sections~\ref{sec:preliminaries}–\ref{sec:objective} detail each
component in turn.

\begin{figure*}[t]
  \centering
  \includegraphics[width=\linewidth]{figs/pipeline_fig.pdf}
  \caption{Overview of the proposed framework. PURSUhInT selects $K$
    hint points offline via $R^2$-CCA clustering of teacher layer
    representations. EVBMF rank selection constructs the CP and Tucker
    students from the teacher's trained weights. Both students are
    co-trained under a shared frozen teacher with BSAT losses at each
    hint point and bidirectional subspace coupling.}
  \label{fig:pipeline}
\end{figure*}


\subsection{Preliminaries: Tensor Decomposition}
\label{sec:preliminaries}
\par
Both student networks are constructed by replacing every eligible
convolution—defined as a $k_H \times k_W$ layer with $k_H > 1$ or
$k_W > 1$ that is not depthwise or grouped—with a sequence of
lightweight operations derived from low-rank tensor factorization.
The factorized architecture is fixed before training begins; only
the layer structure is retained and all weight values are
re-initialized from scratch. This distinguishes the approach from
post-training compression, where decomposition is applied to
pre-trained weights that must then be fine-tuned under a shifting
loss landscape.


\subsubsection{CP Decomposition}
\par
Canonical Polyadic (CP) decomposition \cite{lebedev2014speeding}
approximates a weight tensor
$\mathcal{W} \in \mathbb{R}^{C_{\mathrm{out}} \times C_{\mathrm{in}}
\times k_H \times k_W}$ as a sum of $R$ rank-one terms:
\begin{equation}
    \mathcal{W} \approx \sum_{r=1}^{R}
    \mathbf{a}_r \otimes \mathbf{b}_r \otimes \mathbf{c}_r \otimes \mathbf{d}_r
\end{equation}
where $\otimes$ denotes the outer product and $R$ is the CP rank.
This factorization replaces the original convolution with four
sequential lightweight operations:
\begin{equation}
    \mathbf{x}
    \xrightarrow{1\times 1,\; C_{\mathrm{in}} \to R}
    \xrightarrow{k_H\times 1,\; \mathrm{depthwise}}
    \xrightarrow{1\times k_W,\; \mathrm{depthwise}}
    \xrightarrow{1\times 1,\; R \to C_{\mathrm{out}}}
    \mathbf{y}
\end{equation}
The first and last $1\times 1$ convolutions handle channel mixing;
the two depthwise convolutions apply separable spatial filtering
along each axis independently. The total parameter count reduces
from $C_{\mathrm{out}} \cdot C_{\mathrm{in}} \cdot k_H \cdot k_W$
to $R(C_{\mathrm{in}} + k_H + k_W + C_{\mathrm{out}})$, yielding
compression whenever $R \ll \min(C_{\mathrm{in}}, C_{\mathrm{out}})
\cdot k_H k_W / (k_H + k_W)$. A consequence of the separable
spatial decomposition is that CP strictly limits cross-channel
interactions to the pointwise stages.


\subsubsection{Tucker Decomposition}
\par
Tucker decomposition \cite{kim2015compression} factorizes
$\mathcal{W}$ into a compact core tensor and two orthogonal factor
matrices acting on the channel modes:
\begin{equation}
    \mathcal{W} \approx
    \mathcal{G} \times_1 \mathbf{U}_{\mathrm{out}}
                \times_2 \mathbf{U}_{\mathrm{in}}
\end{equation}
where $\mathcal{G} \in \mathbb{R}^{R_{\mathrm{out}} \times
R_{\mathrm{in}} \times k_H \times k_W}$ is the core,
$\mathbf{U}_{\mathrm{out}} \in \mathbb{R}^{C_{\mathrm{out}} \times
R_{\mathrm{out}}}$ and $\mathbf{U}_{\mathrm{in}} \in
\mathbb{R}^{C_{\mathrm{in}} \times R_{\mathrm{in}}}$ are the factor
matrices. This yields three sequential operations:
\begin{equation}
    \mathbf{x}
    \xrightarrow{1\times 1,\; C_{\mathrm{in}} \to R_{\mathrm{in}}}
    \xrightarrow{k_H \times k_W,\; R_{\mathrm{in}} \to R_{\mathrm{out}}}
    \xrightarrow{1\times 1,\; R_{\mathrm{out}} \to C_{\mathrm{out}}}
    \mathbf{y}
\end{equation}
The full spatial kernel is preserved in the core convolution,
allowing Tucker to maintain cross-channel correlations across
all $k_H \times k_W$ positions—a capacity that CP's separable
decomposition sacrifices. The parameter count reduces from
$C_{\mathrm{out}} \cdot C_{\mathrm{in}} \cdot k_H k_W$ to
$R_{\mathrm{in}} C_{\mathrm{in}} + R_{\mathrm{in}} R_{\mathrm{out}}
k_H k_W + R_{\mathrm{out}} C_{\mathrm{out}}$. The independent
rank pair $(R_{\mathrm{out}}, R_{\mathrm{in}})$ enables separate
control over input and output channel compression, giving Tucker
greater structural flexibility than a single global rank.


\subsection{EVBMF-Guided Rank Selection from the Teacher}
\label{sec:evbmf}
\par
Selecting appropriate decomposition ranks is critical: ranks that are
too low discard representational capacity; ranks that are too high
negate compression gains. Uniform global ratios—setting every layer's
rank to a fixed fraction of its channel count—ignore the heterogeneous
information density across layers. We instead derive per-layer ranks
from the teacher's fully-trained weight spectrum using Empirical
Variational Bayesian Matrix Factorization (EVBMF)
\cite{nakajima2013global} (right branch, Fig.~\ref{fig:pipeline}).

\par
EVBMF treats the weight matrix as a signal-plus-noise observation
under a low-rank prior. Given the mode-0 unfolding
$\mathbf{W} \in \mathbb{R}^{C_{\mathrm{out}} \times (C_{\mathrm{in}}
k_H k_W)}$ of a convolutional weight tensor, EVBMF applies
Bayesian shrinkage to the singular values of $\mathbf{W}$ using the
Marchenko-Pastur noise distribution as a threshold, retaining only
those singular values that exceed the noise floor:
\begin{equation}
    \hat{R} = \bigl|\bigl\{ \sigma_i(\mathbf{W}) :
    \sigma_i(\mathbf{W}) > \sigma^*(\mathbf{W}) \bigr\}\bigr|
\end{equation}
where $\sigma^*(\mathbf{W})$ is the empirical noise threshold
derived from the singular value spectrum of $\mathbf{W}$ and
$\hat{R}$ is the estimated intrinsic rank. This rank captures how
much of the layer's total channel capacity is actively utilized in
the teacher's learned representation.

\par
Because the student and teacher typically differ in depth and width,
student layer $i$ is matched to teacher layer $j$ via positional
interpolation over the set of decomposable convolutions:
\begin{equation}
    j = \mathrm{round}\!\left( i \cdot \frac{T - 1}{S - 1} \right)
    \label{eq:interpolation}
\end{equation}
where $S$ and $T$ are the total numbers of decomposable convolutions
in the student and teacher respectively. The teacher's rank is then
expressed as a utilization fraction of its output channel count and
scaled to the student's channel dimensions:
\begin{equation}
    R_i^{(\mathrm{CP})} = \mathrm{round}\!\left(
        \frac{\hat{R}_j}{C_{\mathrm{out},j}^{(T)}}
        \cdot \max\!\bigl(C_{\mathrm{out},i}^{(S)},
                          C_{\mathrm{in},i}^{(S)}\bigr)
    \right)
\end{equation}
For Tucker decomposition, EVBMF is applied independently to the
mode-0 unfolding (output channels) and the mode-1 unfolding (input
channels) of the teacher layer, yielding a separate rank pair
$(R_{\mathrm{out},i},\, R_{\mathrm{in},i})$ for each student layer.
The resulting rank map is computed once before training begins and
used to construct the decomposed student architectures; the layer
dimensions are fixed for the entire training run. This transfers
the teacher's empirically observed per-layer capacity utilization
directly into the student's structural blueprint, without any
rank adjustment during training.


\subsection{Hint Point Selection via PURSUhInT}
\label{sec:pursuhint}
\par
Rather than selecting distillation targets by hand or applying rigid
spatial heuristics, the proposed framework adopts the PURSUhInT
mechanism \cite{keser2023pursuhint} to identify $K$ non-redundant
hint points from the teacher's layer hierarchy (left branch, Fig.~\ref{fig:pipeline}).
For each of the teacher's $L$ sub-blocks, a forward pass over a
random subset of the training set is performed with spatial average
pooling to obtain an $N \times C_i$ representation matrix
$\mathbf{X}_i$. Pairwise representational distances between all
layers are then computed as $d(i, j) = 1 - R^2_{\mathrm{CCA}}
(\mathbf{X}_i, \mathbf{X}_j)$, where $R^2_{\mathrm{CCA}}$ is the
mean squared Canonical Correlation Analysis coefficient
\cite{kornblith2019similarity} measuring the proportion of variance
explained by the optimal linear mapping between the two
representations. We use K-means clustering to cluster $L$ representations 
using $R^2_{\mathrm{CCA}}$ as the distance function, grouping layers that carry 
redundant semantic information into $K$ clusters. The index closest to the mean index
of each cluster is selected as the centroid, yielding the final hint
point set $\mathcal{H} = \{h_1, \ldots, h_K\} \subset \{1, \ldots,
L\}$. These $K$ indices identify the teacher layers passed to the
BSAT loss at each training step.


\subsection{The BSAT Distillation Loss}
\label{sec:bsat}
\par
The BSAT loss is applied independently at each of the $K$ hint points
$h_k \in \mathcal{H}$. Let $\mathbf{f}_s^k \in \mathbb{R}^{B \times
C_s \times H_s \times W_s}$ and $\mathbf{f}_t^k \in \mathbb{R}^{B
\times C_t \times H_t \times W_t}$ denote the student and teacher
feature maps at hint point $h_k$ for a batch of size $B$. The loss
combines two complementary terms: a spatial attention matching term
that aligns where each network concentrates its discriminative
response, and a batch-subspace alignment term that aligns the
principal directions of variance across the batch—without requiring
$C_s = C_t$ or $H_s = W_s = H_t = W_t$ (Fig.~\ref{fig:bsat}).

\begin{figure*}[t]
  \centering
  \includegraphics[width=\linewidth]{figs/bsat_fig.pdf}
  \caption{Computation of the BSAT loss at a single hint point $k$.
    The AT branch (top) pools the student to the teacher's spatial
    resolution, computes channel-squared attention maps, and minimises
    their MSE. The BSA branch (bottom) independently reshapes each
    feature map into a batch matrix, forms the regularised Gram
    matrix $\hat{\mathbf{G}}^k$, and aligns the energy-adaptive soft
    spectral embeddings $\mathbf{P}^k$ via MSE. The two losses sum to
    $\mathcal{L}_{\mathrm{BSAT}}^k$.}
  \label{fig:bsat}
\end{figure*}

\subsubsection{Spatial Attention Matching}
\par
Following Zagoruyko and Komodakis \cite{zagoruyko2016paying}, the
spatial attention map of a feature tensor is defined by squaring all
activations, averaging across the channel dimension, and
$\ell_2$-normalizing the result (AT band, Fig.~\ref{fig:bsat}). 
When the student and teacher feature maps differ in spatial 
resolution, adaptive average pooling is first
applied to the student to match the teacher's spatial dimensions:
$\tilde{\mathbf{f}}_s^k = \mathrm{AvgPool}\!\left(\mathbf{f}_s^k,
(H_t, W_t)\right)$. The attention maps are then:
\begin{equation}
    \mathbf{a}^k = \frac{\mathbf{v}^k}{\|\mathbf{v}^k\|_2}, \quad
    \mathbf{v}^k = \frac{1}{C}\sum_{c=1}^{C}
    \left(\mathbf{f}^k_{:,c,:,:}\right)^2
    \in \mathbb{R}^{B \times H_t W_t}
\end{equation}
and the attention matching loss at hint point $k$ is:
\begin{equation}
    \mathcal{L}_{\mathrm{AT}}^k =
    \operatorname{MSE}\!\left(\mathbf{a}_s^k,\, \mathbf{a}_t^k\right)
\end{equation}
where $\operatorname{MSE}(\mathbf{A}, \mathbf{B}) =
\frac{1}{|\mathbf{A}|}\|\mathbf{A} - \mathbf{B}\|_F^2$ is the mean
squared error over all elements. This term is sign-invariant by construction
 and requires no
trainable regressor since the spatial dimensions are reconciled by
pooling.

\subsubsection{Batch-Subspace Alignment}
\par
Euclidean feature matching is ill-posed when $C_s \neq C_t$: there
is no canonical correspondence between individual channel activations
of architecturally heterogeneous networks. The batch-subspace
alignment term sidesteps this by working in the batch dimension
rather than the channel dimension (BSA band, Fig.~\ref{fig:bsat}). 
Each feature map is first unfolded into a batch 
matrix $\mathbf{M} \in \mathbb{R}^{B \times D}$ where
$D = C \cdot H \cdot W$, and Frobenius-normalized to remove scale:
\begin{equation}
    \mathbf{M}^k = \frac{\mathbf{f}^k_{[B \times D]}}
    {\|\mathbf{f}^k_{[B \times D]}\|_F}
\end{equation}
The Gram matrix $\mathbf{G}^k = \mathbf{M}^k (\mathbf{M}^k)^\top
\in \mathbb{R}^{B \times B}$ is then formed. Because $\mathbf{G}^k$
is $B \times B$ regardless of $C$, $H$, or $W$, features from
networks with entirely different channel counts and spatial
resolutions are compared in a common space without any trainable
projection. A relative diagonal perturbation is added for numerical stability:
\begin{equation}
    \hat{\mathbf{G}}^k = \mathbf{G}^k + \varepsilon\,\mathbf{I}_B,
    \quad \varepsilon = \max\!\left(10^{-4}\cdot\overline{\operatorname{diag}(\mathbf{G}^k)},\;10^{-6}\right)
\end{equation}
where $\overline{\operatorname{diag}(\mathbf{G}^k)}$ is the mean diagonal of $\mathbf{G}^k$.
Scaling the perturbation with the matrix magnitude provides robust
conditioning against near-degenerate eigenvalue pairs that arise late in
training, where a fixed scalar perturbation would be insufficient.

\par
The symmetric eigendecomposition $\hat{\mathbf{G}}^k = \mathbf{U}\boldsymbol{\Lambda}\mathbf{U}^\top$
(via \texttt{torch.linalg.eigh}, ascending order) yields $B$ eigenpairs.
Rather than retaining a fixed number, the effective rank $q$ is determined
adaptively by an energy threshold $\varepsilon_e \in (0,1]$:
\begin{equation}
    q_{\varepsilon_e} = \min\!\left\{q : \frac{\displaystyle\sum_{i=B-q+1}^{B}\lambda_i}
                                              {\displaystyle\sum_{i=1}^{B}\lambda_i}
                               \geq \varepsilon_e\right\}, \quad
    q = \max\!\left(1,\,\min(q_{\varepsilon_e},\, R,\, B,\, D)\right)
    \label{eq:energy_rank}
\end{equation}
where $R$ is a configurable rank cap. This decouples the effective subspace
dimension from the batch size: near-degenerate features (common in early
training or heavily compressed students) yield a small $q$, while rich
features expand it up to $R$.

\par
The retained top-$q$ eigenvectors $\mathbf{U}_q = \mathbf{U}_{[:,B-q:]}$ are
weighted by their eigenvalue magnitudes via a softmax over the retained
eigenvalues at spectral temperature $\tau > 0$:
\begin{equation}
    \mathbf{w} = q\cdot\operatorname{softmax}\!\!\left(
        \frac{\boldsymbol{\lambda}_{\mathrm{top}\text{-}q}}{\lambda_{\max}\cdot\tau + \delta}
    \right) \in \mathbb{R}^{q}
    \label{eq:soft_weights}
\end{equation}
where $\delta = 10^{-12}$ prevents division by zero. As $\tau\!\to\!\infty$ the
weights become uniform and the representation recovers the standard hard
rank-$R$ projector; as $\tau\!\to\!0$ it collapses to a rank-1 approximation.
The soft spectral embedding of the batch subspace is:
\begin{equation}
    \mathbf{P}^k = \mathbf{U}_q\,\operatorname{diag}(\mathbf{w})\,\mathbf{U}_q^\top
    \in \mathbb{R}^{B \times B}
    \label{eq:proj}
\end{equation}
$\mathbf{P}^k$ is sign-invariant: replacing any column $\mathbf{u}_i$ of
$\mathbf{U}_q$ with $-\mathbf{u}_i$ leaves $\mathbf{P}^k$ unchanged, since
$(-\mathbf{u}_i)(-\mathbf{u}_i)^\top w_i = \mathbf{u}_i\mathbf{u}_i^\top w_i$.
This eliminates the eigenvector sign ambiguity that destabilises direct
eigenvector alignment during backpropagation. The teacher embedding
$\mathbf{P}_t^k$ is computed identically but detached from the computation
graph; gradients flow through the eigendecomposition of $\hat{\mathbf{G}}_s^k$
to the student feature maps. The subspace alignment loss at hint point $k$ is:
\begin{equation}
    \mathcal{L}_{\mathrm{BSA}}^k =
    \operatorname{MSE}\!\left(\mathbf{P}_s^k,\, \mathbf{P}_t^k\right)
\end{equation}

\subsubsection{Combined BSAT Loss}
\par
The total BSAT loss sums both terms across all $K$ hint points (Fig.~\ref{fig:bsat}):
\begin{equation}
    \mathcal{L}_{\mathrm{BSAT}} = \sum_{k=1}^{K}
    \left[ \mathcal{L}_{\mathrm{AT}}^k + \mathcal{L}_{\mathrm{BSA}}^k \right]
\end{equation}
The two terms are complementary: $\mathcal{L}_{\mathrm{AT}}^k$
aligns where the network attends spatially within each sample, while
$\mathcal{L}_{\mathrm{BSA}}^k$ aligns how the batch of samples
organizes itself in the latent space—a structural property that
spatial attention maps do not capture.


\subsection{Heterogeneous Dual-Student Co-Distillation}
\label{sec:dual}

\subsubsection{Dual Architecture Instantiation}
\par
Both student networks are derived from the same base architecture
(e.g., ResNet-18) but decomposed by different methods: one via CP
decomposition ($\mathcal{S}_{\mathrm{CP}}$) and one via Tucker
decomposition ($\mathcal{S}_{\mathrm{Tuck}}$), each using the
per-layer rank map derived in Section~\ref{sec:evbmf}. The two
students are therefore structurally heterogeneous—their layer
configurations, intermediate channel counts, and parameter counts
differ—but their depth and the positions of their hint-point
connectors align with the student hint indices matched to the student's block boundaries. Both are initialized from scratch and trained
simultaneously within a single training loop. They share the same
data loader and teacher forward pass at each step, but maintain
entirely separate optimizers, learning rate schedulers, gradient
scalers, and checkpoint trackers. The two students share a single teacher forward pass per training step, reducing co-training overhead to roughly twice the cost of single-student
distillation rather than the three-pass cost of sequential training. The motivation for simultaneous
training over sequential training is that the bidirectional
coupling term described below creates a live information channel
between the two students that depends on their current-batch
representations; this channel is meaningless if the students are
trained in isolation.

\subsubsection{Bidirectional Batch-Subspace Coupling}
\par
At each hint point $h_k$, the forward passes of both students
produce projectors $\mathbf{P}_{\mathrm{CP}}^k$ and
$\mathbf{P}_{\mathrm{Tuck}}^k$ as a by-product of computing
$\mathcal{L}_{\mathrm{BSAT}}$ (Section~\ref{sec:bsat}, Fig.~\ref{fig:bsat}). 
These projectors encode how each student organizes the current batch in
its latent space. To constrain both students to compress
representations into a shared semantic subspace, a bidirectional
coupling term is added to each student's BSAT loss after both
forward passes are complete:
\begin{equation}
    \mathcal{L}_{\mathrm{coup},\mathrm{CP}} =
    \lambda \sum_{k=1}^{K}
    \operatorname{MSE}\!\left(\mathbf{P}_{\mathrm{CP}}^k,\,
    \overline{\mathbf{P}}_{\mathrm{Tuck}}^k\right)
\end{equation}
\begin{equation}
    \mathcal{L}_{\mathrm{coup},\mathrm{Tuck}} =
    \lambda \sum_{k=1}^{K}
    \operatorname{MSE}\!\left(\mathbf{P}_{\mathrm{Tuck}}^k,\,
    \overline{\mathbf{P}}_{\mathrm{CP}}^k\right)
\end{equation}
where $\overline{\mathbf{P}}$ denotes a stop-gradient operation
(\texttt{detach}) and $\lambda$ is the coupling weight
hyperparameter. The stop-gradient is essential: when minimizing
$\mathcal{L}_{\mathrm{coup},\mathrm{CP}}$, gradients flow through
$\mathbf{P}_{\mathrm{CP}}^k$ to update $\mathcal{S}_{\mathrm{CP}}$,
but $\mathbf{P}_{\mathrm{Tuck}}^k$ is treated as a fixed target
for that step, and vice versa. Neither student is therefore
directly pulled by the other's gradient; each is instead guided
by the other's current subspace as a regression target.

\par
Both coupling terms are computed within the same forward pass
using the current batch's projectors—not projectors from a
previous iteration. This means the coupling targets are always
consistent with the representations the two students produce on
identical input, making the constraint tighter than a momentum
memory bank or lagged target. The two students are not forced to
be identical: $\mathcal{L}_{\mathrm{BSAT}}$ independently aligns
each student with the teacher's subspace, and the coupling term
only requires agreement between the two students' batch
organizations, not between their individual weight matrices or
channel-level activations. Because $\mathbf{P}^k$ is $B \times B$
regardless of decomposition type, the coupling term is
dimension-agnostic and requires no adapter module to bridge the
structural gap between CP and Tucker layers.


\subsection{Training Objective and Optimization}
\label{sec:objective}
\par
Each student is trained with a composite loss combining supervised
classification, output-space distillation, and the BSAT distillation
loss with bidirectional coupling. For student
$\mathcal{S} \in \{\mathcal{S}_{\mathrm{CP}},
\mathcal{S}_{\mathrm{Tuck}}\}$ the full per-step objective is:
\begin{equation}
    \mathcal{L}_{\mathcal{S}} =
    \gamma \,\mathcal{L}_{\mathrm{CE}}
    + \alpha \,\mathcal{L}_{\mathrm{KL}}
    + \beta \!\left[
        \mathcal{L}_{\mathrm{BSAT}}(\mathcal{S})
        + \lambda \sum_{k=1}^{K}
        \operatorname{MSE}\!\left(\mathbf{P}_{\mathcal{S}}^k,\,
        \overline{\mathbf{P}}_{\mathcal{S}'}^k\right)
    \right]
    \label{eq:total_loss}
\end{equation}
where $\mathcal{L}_{\mathrm{CE}}$ is the cross-entropy loss with
ground-truth labels, $\mathcal{L}_{\mathrm{KL}}$ is the
Kullback-Leibler divergence between the student's and teacher's
softened logits at temperature $\tau = 4$, $\mathcal{S}'$ denotes
the peer student, and $\gamma$, $\alpha$, $\beta$, $\lambda$ are
scalar weights. The two students are optimized independently via
separate backward passes; neither shares gradients with the other.

\par
As an alternative to fixed scalar weights, the framework supports
the homoscedastic uncertainty weighting of Kendall
et al.~\cite{kendall2018multi}, which replaces the fixed coefficients
with learnable log-variance parameters $\{s_i\}$:
\begin{equation}
    \mathcal{L}_{\mathcal{S}}^{\mathrm{dyn}} =
    \sum_{i \in \{\mathrm{CE},\,\mathrm{KL},\,\mathrm{AT},\,\mathrm{BSA}\}}
    \left[ e^{-s_i} \mathcal{L}_i + s_i \right]
\end{equation}
where $\mathcal{L}_\mathrm{AT} = \sum_k \mathcal{L}_{\mathrm{AT}}^k$ and
$\mathcal{L}_\mathrm{BSA} = \sum_k \mathcal{L}_{\mathrm{BSA}}^k$.
Each $s_i$ is initialized to zero (effective weight 1) and placed
in a separate parameter group with no weight decay to prevent
divergence. Separating the AT and BSA terms allows the optimizer
to independently resolve their different loss scales. When this
mode is active, $\gamma$, $\alpha$, and $\beta$ are ignored.







%%%%% EXPERIMENTS %%%%%







\section{Experiments}
\label{sec:experiments}

\subsection{Experimental Setup}

\subsubsection{Datasets}
We evaluate the proposed framework on three standard benchmarks.
\textbf{CIFAR-10} and \textbf{CIFAR-100} \cite{krizhevsky2009learning} each
contain 50{,}000 training and 10{,}000 test images at $32{\times}32$
resolution, with 10 and 100 classes respectively.
Standard augmentation is applied to both: random cropping with padding 4
and random horizontal flipping, followed by normalization with per-channel
statistics of the training set.
\textbf{ImageNet-100} is a 100-class subset of ILSVRC-2012
\cite{russakovsky2015imagenet} containing ${\approx}132{,}000$ training and
10{,}000 validation images at $224{\times}224$ resolution.
Training augmentation comprises random resized cropping (scale $[0.1, 1]$,
aspect ratio $[0.8, 1.25]$), random horizontal flipping, and colour jitter
(brightness, contrast, saturation each $\pm 0.4$); validation images are
resized to $256{\times}256$ and centre-cropped to $224{\times}224$.
GPU-accelerated data loading for ImageNet-100 is performed via the NVIDIA
DALI pipeline to eliminate CPU preprocessing bottlenecks.

\subsubsection{Architecture Configurations}
We evaluate six teacher--student architecture pairs on CIFAR-10 and
CIFAR-100, covering both \emph{homogeneous} pairs (teacher and student share
the same architecture family) and \emph{cross-architecture} pairs (different
families), which stress-test the dimension-agnostic property of the BSAT
loss:

\begin{itemize}
  \item \textbf{WRN-40-2 $\to$ WRN-16-2} (homogeneous)
  \item \textbf{WRN-40-2 $\to$ ResNet-8$\times$4} (cross-architecture)
  \item \textbf{ResNet-32$\times$4 $\to$ ResNet-8$\times$4} (homogeneous)
  \item \textbf{ResNet-32$\times$4 $\to$ WRN-16-2} (cross-architecture)
  \item \textbf{ResNet-110 $\to$ ResNet-20} (homogeneous)
  \item \textbf{VGG-13 $\to$ VGG-8} (homogeneous)
\end{itemize}

Since architecture pairs are fixed across datasets, compression metrics
(parameter count, FLOPs, and latency) are identical for CIFAR-10 and
CIFAR-100.
On ImageNet-100, we evaluate two teacher--student configurations:

\begin{itemize}
  \item \textbf{ResNet-34 $\to$ ResNet-18} ($21.34$\,M teacher parameters)
  \item \textbf{ResNet-50 $\to$ ResNet-18} ($23.71$\,M teacher parameters)
\end{itemize}

In all cases, each student skeleton is instantiated twice: once as a
CP-decomposed network $\mathcal{S}_{\mathrm{CP}}$ and once as a
Tucker-decomposed network $\mathcal{S}_{\mathrm{Tuck}}$, with per-layer ranks
derived from the respective teacher's trained weight spectrum via EVBMF
(Section~\ref{sec:evbmf}). All student weights are initialised from scratch;
no pre-trained weights are transferred.

\subsubsection{Training Protocol}
Both students are trained simultaneously with SGD (momentum $0.9$).
On \textbf{CIFAR-10} and \textbf{CIFAR-100}, the learning rate is $0.05$
with weight decay $5\!\times\!10^{-4}$, decayed by $\times\!0.1$ at epochs
150, 180, and 210 over 240 total epochs with a 5-epoch linear warmup.
Batch size is 64 on a single GPU.
On \textbf{ImageNet-100}, the base learning rate is $0.1$, linearly scaled
as $\eta_0\!\cdot\!\text{batch\_size}\!\times\!\text{world\_size}/256$, with
a 5-epoch linear warmup followed by $\times\!0.1$ decay at epochs 30, 60,
80, and 90 over 100 total epochs, with weight decay $10^{-4}$ and batch
size 256 per GPU.
Distributed training uses PyTorch DDP with NCCL.
Mixed-precision (AMP) training with per-student gradient scalers is applied
throughout.

\subsubsection{Distillation Hyperparameters}
The KD temperature is $\tau_\mathrm{KD} = 4$, the BSAT subspace rank cap is
$R = 8$, the energy threshold is $\varepsilon_e = 0.9$, and the spectral
temperature is $\tau = 1.0$ (Eqs.~\ref{eq:energy_rank}--\ref{eq:soft_weights}).
Rather than tuning fixed loss coefficients $(\gamma, \alpha, \beta)$, we
apply the homoscedastic uncertainty weighting of Kendall
et al.~\cite{kendall2018multi} with four separate log-variance parameters
($s_\mathrm{CE}$, $s_\mathrm{KL}$, $s_\mathrm{AT}$, $s_\mathrm{BSA}$) throughout.
The bidirectional coupling weight is fixed at $\lambda = 1.0$
(Eq.~\ref{eq:total_loss}).

\subsubsection{Hint Point Selection}
PURSUhInT is applied offline to each trained teacher using $R^2_{\mathrm{CCA}}$
as the representational distance metric with $K = 3$ clusters for CIFAR and
$K = 4$ clusters for ImageNet-100.
Since hint point selection depends only on the teacher architecture and not
on the dataset, the same hint indices are reused for CIFAR-10 and CIFAR-100
under a shared teacher.
For the WRN-40-2 teacher the selected indices are $\{5, 10, 16\}$;
for ResNet-32$\times$4 they are $\{5, 8, 13\}$;
for ResNet-110 they are $\{1, 6, 14\}$.
For the ResNet-34 and ResNet-50 ImageNet teachers the selected indices
are $\{3, 8, 11, 16\}$ and $\{8, 11, 12, 16\}$ respectively.

\subsubsection{Baselines}
We compare the proposed BSAT loss against the following distillation
objectives applied to identically compressed CP and Tucker students:
(i)~\textbf{KD} \cite{hinton2015distilling}: softened-logit KL divergence
only (no intermediate supervision);
(ii)~\textbf{FitNet} \cite{romero2015fitnetshintsdeepnets}: channel-wise
feature regression with $1{\times}1$ convolutional adapters at each hint
point;
(iii)~\textbf{AT} \cite{zagoruyko2016paying}: spatial attention map MSE at
each hint point, without the BSA term;
(iv)~\textbf{VID} \cite{ahn2019variational}: variational information
distillation;
(v)~\textbf{WSL} \cite{zhou2021rethinking}: weighted soft labels that
dynamically calibrate loss weights by sample difficulty.
All baselines train each student independently (no bidirectional coupling
term). Accuracy, parameter count, FLOPs, and latency are measured on the
same hardware with the same evaluation pipeline for all methods.

%% ─────────────────────────────────────────────────────────────────────────
\subsection{Main Results on ImageNet-100}
\label{sec:results_imagenet}

Table~\ref{tab:imagenet100} reports results for both ImageNet-100
configurations. Inference latency is measured on a single CPU core with
batch size 1, averaged over 50 runs after 10 warmup iterations.

\begin{table*}[t]
  \centering
  \caption{Results on ImageNet-100. Params and FLOPs for all student methods
    reflect the \emph{deployed} model (no adapter overhead).
    Compression ratios are relative to the respective teacher.}
  \label{tab:imagenet100}
  \setlength{\tabcolsep}{4.5pt}
  \begin{tabular}{llcccccc}
    \toprule
    \textbf{Method} & \textbf{Decomp.} &
    \textbf{Top-1 (\%)} & \textbf{Top-5 (\%)} &
    \textbf{Params} & \textbf{FLOPs} &
    \textbf{Params Compr.} & \textbf{FLOPs Compr.} \\
    \midrule
    \multicolumn{8}{l}{\textit{ResNet-34 $\to$ ResNet-18}} \\
    \midrule
    Teacher (ResNet-34) & —
      & 81.72 & 95.38 & 21.34\,M & 3.68\,G & $1.0\times$ & $1.0\times$ \\
    \cline{1-8}
    KD \cite{hinton2015distilling}                       & \multirow{6}{*}{CP}     & 77.16 & 93.68 & \multirow{6}{*}{1.33\,M} & \multirow{6}{*}{277.99\,M} & \multirow{6}{*}{$16.04\times$} & \multirow{6}{*}{$13.23\times$} \\
    FitNet \cite{romero2015fitnetshintsdeepnets}         &                         & 78.48 & 94.70 & & & & \\
    AT \cite{zagoruyko2016paying}                        &                         & 79.50 & 95.00 & & & & \\
    VID \cite{ahn2019variational}                        &                         & 78.52 & 94.54 & & & & \\
    WSL \cite{zhou2021rethinking}                        &                         & 77.44 & 94.42 & & & & \\
    \textbf{Ours (BSAT)}                                 &                         & 80.10 & 94.90 & & & & \\
    \cline{1-8}
    KD \cite{hinton2015distilling}                       & \multirow{6}{*}{Tucker} & 77.26 & 94.14 & \multirow{6}{*}{3.25\,M} & \multirow{6}{*}{655.50\,M} & \multirow{6}{*}{$6.56\times$} & \multirow{6}{*}{$5.60\times$} \\
    FitNet \cite{romero2015fitnetshintsdeepnets}         &                         & 80.22 & 95.34 & & & & \\
    AT \cite{zagoruyko2016paying}                        &                         & 80.12 & 95.32 & & & & \\
    VID \cite{ahn2019variational}                        &                         & 79.88 & 95.14 & & & & \\
    WSL \cite{zhou2021rethinking}                        &                         & 79.44 & 95.16 & & & & \\
    \textbf{Ours (BSAT)}                                 &                         & 80.80 & 95.34 & & & & \\
    \midrule
    \multicolumn{8}{l}{\textit{ResNet-50 $\to$ ResNet-18}} \\
    \midrule
    Teacher (ResNet-50) & —
      & 82.54 & 95.96 & 23.71\,M & 4.13\,G & $1.0\times$ & $1.0\times$ \\
    \cline{1-8}
    KD \cite{hinton2015distilling}                       & \multirow{6}{*}{CP}     & 78.10 & 93.94 & \multirow{6}{*}{1.35\,M} & \multirow{6}{*}{270.15\,M} & \multirow{6}{*}{$17.59\times$} & \multirow{6}{*}{$15.29\times$} \\
    FitNet \cite{romero2015fitnetshintsdeepnets}         &                         & 78.48 & 94.44 & & & & \\
    AT \cite{zagoruyko2016paying}                        &                         & 79.60 & 94.84 & & & & \\
    VID \cite{ahn2019variational}                        &                         & 78.14 & 94.04 & & & & \\
    WSL \cite{zhou2021rethinking}                        &                         & 78.58 & 94.52 & & & & \\
    \textbf{Ours (BSAT)}                                 &                         & 79.34 & 95.12 & & & & \\
    \cline{1-8}
    KD \cite{hinton2015distilling}                       & \multirow{6}{*}{Tucker} & 77.66 & 93.94 & \multirow{6}{*}{3.18\,M} & \multirow{6}{*}{600.63\,M} & \multirow{6}{*}{$7.46\times$} & \multirow{6}{*}{$6.88\times$} \\
    FitNet \cite{romero2015fitnetshintsdeepnets}         &                         & 80.10 & 95.28 & & & & \\
    AT \cite{zagoruyko2016paying}                        &                         & 81.04 & 95.56 & & & & \\
    VID \cite{ahn2019variational}                        &                         & 80.18 & 95.48 & & & & \\
    WSL \cite{zhou2021rethinking}                        &                         & 80.24 & 95.34 & & & & \\
    \textbf{Ours (BSAT)}                                 &                         & 80.94 & 95.74 & & & & \\
    \bottomrule
  \end{tabular}
\end{table*}


For the ResNet-34 $\to$ ResNet-18 configuration, the CP-decomposed student
reaches \textbf{80.10\%} Top-1—within 1.62\,pp of the teacher (81.72\%)—while
compressing parameters by $16.0\times$ (1.33\,M vs.\ 21.34\,M) and FLOPs by
$13.2\times$ (278\,M vs.\ 3.68\,G), reducing CPU latency from 265.98\,ms to
135.88\,ms.
The Tucker-decomposed student reaches \textbf{80.80\%} Top-1 at a more
moderate $6.6\times$ parameter compression, reflecting Tucker's retention of
cross-channel correlations at the cost of a larger core tensor.
Crucially, neither student requires any trainable adapter module; the
dimension mismatch between teacher and students is resolved entirely by the
batch-mode Gram formulation of the BSA term.

The Tucker student's accuracy advantage over the CP student ($+0.70$\,pp)
mirrors the architectural trade-off discussed in
Section~\ref{sec:preliminaries}: the dense $k_H{\times}k_W$ core convolution
in Tucker preserves cross-channel structure that CP's separable decomposition
discards.
At convergence, the learned dynamic loss weights reveal that intermediate
feature supervision dominates: the BSA effective weight
($e^{-s_{\mathrm{BSA}}}$) far exceeds the KL divergence weight
($e^{-s_{\mathrm{KL}}}$), while the cross-entropy weight stabilises at a
moderate value—consistent with output-space distillation serving as a coarse
regulariser once the students develop structured representations.

%% ─────────────────────────────────────────────────────────────────────────
\subsection{Main Results on CIFAR-10 and CIFAR-100}
\label{sec:results_cifar}

Table~\ref{tab:cifar} reports results for all six teacher--student
configurations on CIFAR-10 and CIFAR-100. Because the same architecture
pairs are used across both datasets, compression metrics (Params, FLOPs,
and latency) are identical for the two datasets and are listed only once
per pair.

%% ════════════════════════════════════════════════════════════════════════
%%  TABLE 1 — CIFAR-100
%% ════════════════════════════════════════════════════════════════════════
%% ── Table 1 of 2 ────────────────────────────────────────────────────────────
\begin{table*}[t]
  \centering
  \caption{Results on CIFAR-100 (configurations 1--3 of 6).
    Compression metrics derived from VBMF ranks per run.
    ``—'' denotes results pending experimental runs.}
  \label{tab:cifar100_a}
  \setlength{\tabcolsep}{4pt}
  \begin{tabular}{llcccccc}
    \toprule
    \textbf{Method} & \textbf{Decomp.} &
    \textbf{Top-1 (\%)} & \textbf{Top-5 (\%)} &
    \textbf{Params} & \textbf{FLOPs} &
    \textbf{Params Compr.} & \textbf{FLOPs Compr.} \\
    \midrule

    %% ── WRN-40-2 → WRN-16-2 ──────────────────────────────────────────────
    \multicolumn{8}{l}{\textit{WRN-40-2 $\to$ WRN-16-2}} \\
    \midrule
    Teacher (WRN-40-2)   & — & 76.79 & 93.60 & 2.26\,M & 330.43\,M & $1.00\times$ & $1.00\times$ \\
    \cline{1-8}
    KD \cite{hinton2015distilling}                 & \multirow{6}{*}{CP}     & 69.01 & 91.63 & \multirow{6}{*}{111.80\,K} & \multirow{6}{*}{17.20\,M} & \multirow{6}{*}{$20.27\times$} & \multirow{6}{*}{$19.21\times$} \\
    FitNet \cite{romero2015fitnetshintsdeepnets}   &                         & 71.47 & 92.80 & & & & \\
    AT \cite{zagoruyko2016paying}                  &                         & 71.85 & 93.00 & & & & \\
    VID \cite{ahn2019variational}                  &                         & 71.34 & 92.47 & & & & \\
    WSL \cite{zhou2021rethinking}                  &                         & 70.59 & 92.79 & & & & \\
    \textbf{Ours (BSAT)}                           &                         & — & — & & & & \\
    \cline{1-8}
    KD \cite{hinton2015distilling}                 & \multirow{6}{*}{Tucker} & 68.79 & 91.17 & \multirow{6}{*}{332.76\,K} & \multirow{6}{*}{52.03\,M} & \multirow{6}{*}{$6.78\times$} & \multirow{6}{*}{$6.35\times$} \\
    FitNet \cite{romero2015fitnetshintsdeepnets}   &                         & 71.95 & 92.97 & & & & \\
    AT \cite{zagoruyko2016paying}                  &                         & 72.23 & 93.16 & & & & \\
    VID \cite{ahn2019variational}                  &                         & 71.33 & 92.75 & & & & \\
    WSL \cite{zhou2021rethinking}                  &                         & 71.82 & 92.87 & & & & \\
    \textbf{Ours (BSAT)}                           &                         & — & — & & & & \\
    \midrule

    %% ── WRN-40-2 → ResNet-8×4 ────────────────────────────────────────────
    \multicolumn{8}{l}{\textit{WRN-40-2 $\to$ ResNet-8$\times$4 {\scriptsize(cross-arch.)}}} \\
    \midrule
    Teacher (WRN-40-2)   & — & 76.01 & 93.79 & 2.26\,M & 330.43\,M & $1.00\times$ & $1.00\times$ \\
    \cline{1-8}
    KD \cite{hinton2015distilling}                 & \multirow{6}{*}{CP}     & 73.27 & 93.45 & \multirow{6}{*}{258.77\,K} & \multirow{6}{*}{39.05\,M} & \multirow{6}{*}{$8.71\times$} & \multirow{6}{*}{$8.46\times$} \\
    FitNet \cite{romero2015fitnetshintsdeepnets}   &                         & 72.65 & 93.11 & & & & \\
    AT \cite{zagoruyko2016paying}                  &                         & 73.56 & 93.50 & & & & \\
    VID \cite{ahn2019variational}                  &                         & 72.29 & 93.08 & & & & \\
    WSL \cite{zhou2021rethinking}                  &                         & 72.30 & 93.11 & & & & \\
    \textbf{Ours (BSAT)}                           &                         & — & — & & & & \\
    \cline{1-8}
    KD \cite{hinton2015distilling}                 & \multirow{6}{*}{Tucker} & 71.61 & 93.13 & \multirow{6}{*}{726.67\,K} & \multirow{6}{*}{92.24\,M} & \multirow{6}{*}{$3.10\times$} & \multirow{6}{*}{$3.58\times$} \\
    FitNet \cite{romero2015fitnetshintsdeepnets}   &                         & 72.03 & 92.82 & & & & \\
    AT \cite{zagoruyko2016paying}                  &                         & 72.20 & 93.28 & & & & \\
    VID \cite{ahn2019variational}                  &                         & 72.25 & 92.98 & & & & \\
    WSL \cite{zhou2021rethinking}                  &                         & 71.22 & 92.81 & & & & \\
    \textbf{Ours (BSAT)}                           &                         & — & — & & & & \\
    \midrule

    %% ── ResNet-32×4 → ResNet-8×4 ─────────────────────────────────────────
    \multicolumn{8}{l}{\textit{ResNet-32$\times$4 $\to$ ResNet-8$\times$4}} \\
    \midrule
    Teacher (ResNet-32$\times$4) & — & 79.61 & 94.94 & 7.43\,M & 1.09\,G & $1.00\times$ & $1.00\times$ \\
    \cline{1-8}
    KD \cite{hinton2015distilling}                 & \multirow{6}{*}{CP}     & 71.88 & 92.40 & \multirow{6}{*}{252.53\,K} & \multirow{6}{*}{48.77\,M} & \multirow{6}{*}{$29.44\times$}    & \multirow{6}{*}{$22.31\times$} \\
    FitNet \cite{romero2015fitnetshintsdeepnets}   &                         & 72.44 & 92.78 & & & & \\
    AT \cite{zagoruyko2016paying}                  &                         & 73.54 & 93.49 & & & & \\
    VID \cite{ahn2019variational}                  &                         & 72.79 & 92.92 & & & & \\
    WSL \cite{zhou2021rethinking}                  &                         & 73.08 & 93.37 & & & & \\
    \textbf{Ours (BSAT)}                           &                         & — & — & & & & \\
    \cline{1-8}
    KD \cite{hinton2015distilling}                 & \multirow{6}{*}{Tucker} & 70.85 & 92.59 & \multirow{6}{*}{679.70\,K} & \multirow{6}{*}{130.57\,M} & \multirow{6}{*}{$10.94\times$}   & \multirow{6}{*}{$8.33\times$} \\
    FitNet \cite{romero2015fitnetshintsdeepnets}   &                         & 71.70 & 92.98 & & & & \\
    AT \cite{zagoruyko2016paying}                  &                         & 73.14 & 93.63 & & & & \\
    VID \cite{ahn2019variational}                  &                         & 71.44 & 92.54 & & & & \\
    WSL \cite{zhou2021rethinking}                  &                         & 72.17 & 93.22 & & & & \\
    \textbf{Ours (BSAT)}                           &                         & — & — & & & & \\
    \bottomrule
  \end{tabular}
\end{table*}

%% ── Table 2 of 2 ────────────────────────────────────────────────────────────
\begin{table*}[t]
  \centering
  \caption{Results on CIFAR-100 (configurations 4--6 of 6).
    Compression metrics derived from VBMF ranks per run.
    ``—'' denotes results pending experimental runs.}
  \label{tab:cifar100_b}
  \setlength{\tabcolsep}{4pt}
  \begin{tabular}{llcccccc}
    \toprule
    \textbf{Method} & \textbf{Decomp.} &
    \textbf{Top-1 (\%)} & \textbf{Top-5 (\%)} &
    \textbf{Params} & \textbf{FLOPs} &
    \textbf{Params Compr.} & \textbf{FLOPs Compr.} \\
    \midrule

    %% ── ResNet-32×4 → WRN-16-2 ───────────────────────────────────────────
    \multicolumn{8}{l}{\textit{ResNet-32$\times$4 $\to$ WRN-16-2 {\scriptsize(cross-arch.)}}} \\
    \midrule
    Teacher (ResNet-32$\times$4) & — & 79.64 & 94.69 & 7.43\,M & 1.09\,G & $1.00\times$ & $1.00\times$ \\
    \cline{1-8}
    KD \cite{hinton2015distilling}                 & \multirow{6}{*}{CP}     & 66.02 & 89.78 & \multirow{6}{*}{116.09\,K} & \multirow{6}{*}{20.83\,M} & \multirow{6}{*}{$64.04\times$}    & \multirow{6}{*}{$52.24\times$} \\
    FitNet \cite{romero2015fitnetshintsdeepnets}   &                         & 71.36 & 92.50 & & & & \\
    AT \cite{zagoruyko2016paying}                  &                         & 72.01 & 93.26 & & & & \\
    VID \cite{ahn2019variational}                  &                         & 69.29 & 91.46 & & & & \\
    WSL \cite{zhou2021rethinking}                  &                         & 71.47 & 92.71 & & & & \\
    \textbf{Ours (BSAT)}                           &                         & — & — & & & & \\
    \cline{1-8}
    KD \cite{hinton2015distilling}                 & \multirow{6}{*}{Tucker} & 68.25 & 91.23 & \multirow{6}{*}{340.19\,K} & \multirow{6}{*}{70.65\,M} & \multirow{6}{*}{$21.85\times$}    & \multirow{6}{*}{$15.40\times$} \\
    FitNet \cite{romero2015fitnetshintsdeepnets}   &                         & 71.59 & 92.92 & & & & \\
    AT \cite{zagoruyko2016paying}                  &                         & 73.06 & 93.46 & & & & \\
    VID \cite{ahn2019variational}                  &                         & 68.52 & 91.58 & & & & \\
    WSL \cite{zhou2021rethinking}                  &                         & 72.05 & 93.06 & & & & \\
    \textbf{Ours (BSAT)}                           &                         & — & — & & & & \\
    \midrule

    %% ── ResNet-110 → ResNet-20 ───────────────────────────────────────────
    \multicolumn{8}{l}{\textit{ResNet-110 $\to$ ResNet-20}} \\
    \midrule
    Teacher (ResNet-110) & — & 74.30 & 93.21 & 1.74\,M & 257.40\,M & $1.00\times$ & $1.00\times$ \\
    \cline{1-8}
    KD \cite{hinton2015distilling}                 & \multirow{6}{*}{CP}     & 63.64 & 89.00 & \multirow{6}{*}{47.60\,K} & \multirow{6}{*}{8.35\,M} & \multirow{6}{*}{$36.48\times$}     & \multirow{6}{*}{$30.83\times$} \\
    FitNet \cite{romero2015fitnetshintsdeepnets}   &                         & 66.67 & 90.60 & & & & \\
    AT \cite{zagoruyko2016paying}                  &                         & 66.06 & 90.37 & & & & \\
    VID \cite{ahn2019variational}                  &                         & 66.74 & 91.12 & & & & \\
    WSL \cite{zhou2021rethinking}                  &                         & 65.23 & 90.25 & & & & \\
    \textbf{Ours (BSAT)}                           &                         & — & — & & & & \\
    \cline{1-8}
    KD \cite{hinton2015distilling}                 & \multirow{6}{*}{Tucker} & 65.79 & 90.61 & \multirow{6}{*}{127.24\,K} & \multirow{6}{*}{23.00\,M} & \multirow{6}{*}{$13.65\times$}    & \multirow{6}{*}{$11.19\times$} \\
    FitNet \cite{romero2015fitnetshintsdeepnets}   &                         & 67.43 & 91.35 & & & & \\
    AT \cite{zagoruyko2016paying}                  &                         & 67.15 & 90.92 & & & & \\
    VID \cite{ahn2019variational}                  &                         & 67.09 & 91.33 & & & & \\
    WSL \cite{zhou2021rethinking}                  &                         & 66.72 & 90.72 & & & & \\
    \textbf{Ours (BSAT)}                           &                         & — & — & & & & \\
    \midrule

    %% ── VGG-13 → VGG-8 ──────────────────────────────────────────────────
    \multicolumn{8}{l}{\textit{VGG-13 $\to$ VGG-8}} \\
    \midrule
    Teacher (VGG-13)     & — & 75.21 & 92.49 & 9.46\,M & 285.99\,M & $1.00\times$ & $1.00\times$ \\
    \cline{1-8}
    KD \cite{hinton2015distilling}                 & \multirow{6}{*}{CP}     & 70.79 & 91.44 & \multirow{6}{*}{574.08\,K} & \multirow{6}{*}{14.79\,M} & \multirow{6}{*}{$16.48\times$}    & \multirow{6}{*}{$19.34\times$} \\
    FitNet \cite{romero2015fitnetshintsdeepnets}   &                         & 70.20 & 91.34 & & & & \\
    AT \cite{zagoruyko2016paying}                  &                         & 70.26 & 91.49 & & & & \\
    VID \cite{ahn2019variational}                  &                         & 70.28 & 91.35 & & & & \\
    WSL \cite{zhou2021rethinking}                  &                         & 69.10 & 90.97 & & & & \\
    \textbf{Ours (BSAT)}                           &                         & — & — & & & & \\
    \cline{1-8}
    KD \cite{hinton2015distilling}                 & \multirow{6}{*}{Tucker} & 69.85 & 91.29 & \multirow{6}{*}{1.63\,M} & \multirow{6}{*}{40.21\,M} & \multirow{6}{*}{$5.82\times$}     & \multirow{6}{*}{$7.11\times$} \\
    FitNet \cite{romero2015fitnetshintsdeepnets}   &                         & 68.94 & 90.60 & & & & \\
    AT \cite{zagoruyko2016paying}                  &                         & 69.22 & 91.03 & & & & \\
    VID \cite{ahn2019variational}                  &                         & 68.55 & 90.86 & & & & \\
    WSL \cite{zhou2021rethinking}                  &                         & 67.98 & 90.36 & & & & \\
    \textbf{Ours (BSAT)}                           &                         & — & — & & & & \\
    \bottomrule
  \end{tabular}
\end{table*}






Across the six CIFAR configurations, the EVBMF rank estimator selects very
low intrinsic ranks for $32{\times}32$ feature maps, yielding parameter
compressions between $65\times$ and $106\times$ for same-architecture pairs.
These extreme compression ratios make CIFAR an important stress test for the
distillation loss: at $100\times$ compression, logit-level distillation alone
is insufficient, and hint-point supervision becomes the primary mechanism
for knowledge transfer.

The two cross-architecture configurations—WRN-40-2 $\to$ ResNet-8$\times$4
and ResNet-32$\times$4 $\to$ WRN-16-2—demonstrate a key property of the
BSAT loss: because both the AT and BSA terms operate on the batch dimension
($B \times H_tW_t$ and $B \times B$ respectively) rather than the channel
dimension, they impose no structural requirement on the student architecture.
The cross-architecture student receives the same EVBMF-guided rank assignment
(via positional interpolation) and is trained with the same objective as the
homogeneous students, with no modification to the loss or training loop.

%% ─────────────────────────────────────────────────────────────────────────
\subsection{Ablation Studies}
\label{sec:ablation}

We ablate the main design choices on the ImageNet-100 ResNet-34 $\to$
ResNet-18 configuration. All variants are trained with the same optimisation
protocol; only the distillation objective or rank selection method is
modified.

\subsubsection{Contribution of Each Loss Component}

Table~\ref{tab:ablation_loss} isolates the two terms of the BSAT loss and the
bidirectional coupling term. ``AT-only'' removes the BSA term and trains each
student independently; ``BSA-only'' removes the spatial attention term;
``No coupling'' uses the full BSAT loss but sets $\lambda = 0$.

\begin{table}[t]
  \centering
  \caption{Ablation of BSAT loss components on ImageNet-100
    (ResNet-34 $\to$ ResNet-18). Top-1 accuracy (\%).}
  \label{tab:ablation_loss}
  \begin{tabular}{lcc}
    \toprule
    \textbf{Variant} & \textbf{CP Top-1} & \textbf{Tucker Top-1} \\
    \midrule
    KD only (no hint supervision)         & — & — \\
    AT only ($\mathcal{L}_{\mathrm{AT}}$, no BSA)  & — & — \\
    BSA only ($\mathcal{L}_{\mathrm{BSA}}$, no AT) & — & — \\
    BSAT, no coupling ($\lambda=0$)        & — & — \\
    \textbf{BSAT + coupling (ours)}        & \textbf{80.10} & \textbf{80.80} \\
    \bottomrule
  \end{tabular}
\end{table}

The KD-only baseline establishes the floor: without intermediate feature
supervision, heavily factorised students cannot overcome the capacity gap
through logit matching alone. Introducing AT provides coarse spatial guidance
but cannot resolve channel-dimension mismatches without additional adapters;
the BSA term bypasses this entirely via the batch-mode Gram formulation.
Removing the coupling term degrades both students, with Tucker benefiting
more from the CP student's complementary factorisation signal—consistent
with Tucker's richer parameterisation having more capacity to absorb a
structural prior.

\subsubsection{EVBMF vs.\ Uniform Rank Selection}

Table~\ref{tab:ablation_rank} compares EVBMF-guided per-layer rank selection
against two uniform alternatives: a fixed fraction of the maximum channel
count (``Uniform-0.5'') and a conservative fixed fraction (``Uniform-0.25'').

\begin{table}[t]
  \centering
  \caption{Rank selection strategy ablation on ImageNet-100
    (ResNet-34 $\to$ ResNet-18). Top-1 accuracy (\%).}
  \label{tab:ablation_rank}
  \begin{tabular}{lccc}
    \toprule
    \textbf{Rank Strategy} & \textbf{Params} & \textbf{CP Top-1} & \textbf{Tucker Top-1} \\
    \midrule
    Uniform-0.25          & —                & — & — \\
    Uniform-0.5           & —                & — & — \\
    \textbf{EVBMF (ours)} & 1.33\,M / 3.25\,M & \textbf{80.10} & \textbf{80.80} \\
    \bottomrule
  \end{tabular}
\end{table}

EVBMF-derived ranks improve accuracy over uniform ratios by reflecting the
teacher's heterogeneous per-layer capacity utilisation: early layers receive
higher ranks (preserving fine-grained spatial structure) while intermediate
layers receive significantly fewer rank units. Uniform allocation wastes rank
budget in redundant layers and under-provisions layers with high intrinsic
dimensionality.

\subsubsection{Number of Hint Points $K$}

Table~\ref{tab:ablation_hints} varies the number of PURSUhInT hint points
$K \in \{2, 3, 4, 5\}$ on the ImageNet-100 ResNet-34 $\to$ ResNet-18
configuration.

\begin{table}[t]
  \centering
  \caption{Effect of hint point count $K$ on ImageNet-100 Top-1 accuracy (\%).}
  \label{tab:ablation_hints}
  \begin{tabular}{ccc}
    \toprule
    $K$ & \textbf{CP Top-1} & \textbf{Tucker Top-1} \\
    \midrule
    2              & — & — \\
    3              & — & — \\
    \textbf{4 (ours)} & \textbf{80.10} & \textbf{80.80} \\
    5              & — & — \\
    \bottomrule
  \end{tabular}
\end{table}

Performance peaks at $K\!=\!4$: once all semantically distinct teacher layers
have been covered by a cluster centroid, additional hint points sample from
within existing clusters, introducing redundant supervision without adding
structural information.

\subsubsection{Dynamic vs.\ Fixed Loss Weighting}

Replacing Kendall et al.\ uncertainty weighting with fixed coefficients on
ImageNet-100 yields a Top-1 drop of $-$\,pp (CP) and $-$\,pp (Tucker).
The four dynamic weights ($s_\mathrm{CE}$, $s_\mathrm{KL}$, $s_\mathrm{AT}$,
$s_\mathrm{BSA}$) converge to effective values reported in the final
experimental run (see supplementary logs). The BSA effective weight
$e^{-s_\mathrm{BSA}}$ substantially exceeds the KL weight $e^{-s_\mathrm{KL}}$,
reflecting the loss-scale imbalance between cross-entropy and subspace MSE
under high-rank compression; dynamic weighting resolves this automatically
without exhaustive grid search.

%% ─────────────────────────────────────────────────────────────────────────
\subsection{Efficiency Analysis}
\label{sec:efficiency}

\subsubsection{Compression--Accuracy Trade-off}
Fig.~\ref{fig:scatter} plots Top-1 accuracy against parameter compression
ratio for all methods on ImageNet-100 (ResNet-34 $\to$ ResNet-18). The CP
student occupies the $>\!10\times$ parameter regime and the Tucker student
the $4$--$8\times$ regime; together they trace two points on the Pareto
frontier, offering a practitioner the choice between extreme parameter
reduction (CP, $16\times$, 80.10\%) and moderate compression with higher
fidelity (Tucker, $6.6\times$, 80.80\%).

\subsubsection{Per-Layer Rank Distribution}
Fig.~\ref{fig:ranks} visualises the EVBMF-estimated intrinsic rank for each
decomposable layer of the ResNet-34 teacher projected onto the ResNet-18
student via positional interpolation (Eq.~\ref{eq:interpolation}).
Early layers receive relatively high rank estimates (high spatial-frequency
content), while intermediate layers exhibit the steepest rank collapse
(low intrinsic dimensionality of mid-level semantic representations).
Relative to a uniform $50\%$ ratio, EVBMF allocates substantially fewer
rank units to layers 5--12 while preserving capacity in layers 13--16.

\subsubsection{Training Convergence}
Fig.~\ref{fig:convergence} shows Top-1 validation accuracy curves for both
students over 100 training epochs on ImageNet-100. The CP student converges
slightly faster in early epochs, while the Tucker student overtakes it by
approximately epoch 35 and maintains a consistent $+0.6$\,pp advantage
through convergence. The gap to the teacher (81.72\%) stabilises after
epoch 60, indicating that the remaining deficit is determined by the
compression ratio rather than training duration.

\subsubsection{Inference Efficiency}
The CP student reduces single-image CPU latency from 265.98\,ms (teacher)
to 135.88\,ms—a $1.96\times$ speedup—while compressing parameters by
$16.0\times$ and FLOPs by $13.2\times$. The Tucker student achieves a
slightly larger latency reduction (125.91\,ms, $2.11\times$ speedup) despite
its larger parameter count, because the Tucker-factorised convolution sequence
is more cache-friendly on CPU. Both students eliminate the need for any
runtime adapter module; the deployed model is the bare decomposed backbone
with no training-time scaffolding.












\section{Conclusion}
\label{sec:conclusion}

We presented a heterogeneous structural co-distillation framework that jointly
compresses a pre-trained teacher into a CP-decomposed and a Tucker-decomposed
student without pre-trained weight initialisation, trainable feature regressors,
or shared channel dimensions. The proposed BSAT loss aligns sign-invariant soft
spectral embeddings of the batch Gram matrix at PURSUhInT-selected hint points,
accommodating arbitrary channel and spatial mismatches through a purely
batch-mode formulation. Energy-adaptive rank selection decouples the effective
subspace dimension from the batch size, while eigenvalue-proportional soft
weighting provides gradient-smooth alignment. A bidirectional coupling term
constrains both students to organise their representations into a consistent
semantic subspace, exploiting the complementary structural properties of CP and
Tucker factorisation within a single training run. Per-layer ranks derived from
the teacher's weight spectrum via EVBMF transfer the teacher's empirically
observed capacity utilisation into each student's structural blueprint,
consistently outperforming fixed global rank ratios.

On ImageNet-100, the CP student achieves 80.10\% Top-1 accuracy at $16\times$
parameter compression and $13.2\times$ FLOPs reduction relative to the
ResNet-34 teacher ($1.96\times$ CPU latency speedup), while the Tucker student
reaches 80.80\% at a more moderate $6.6\times$ compression ($2.11\times$
speedup). Both results are obtained with no runtime adapter module; the deployed
model is the bare decomposed backbone.

Several directions remain open for future work. The $B\!\times\!B$ Gram
formulation scales quadratically with batch size; mini-batch Gram sketching
would reduce memory for large-batch regimes. Extending the framework to
transformer architectures requires adapting the hint-point extractor to
attention layers, where the spatial dimension is replaced by sequence position.
Combining tensor decomposition with post-training quantisation or structured
pruning could push compression ratios beyond those achievable by structural
factorisation alone.

\bibliographystyle{IEEEtran}
\bibliography{references}

\end{document}